\section{Merge depth replaces schedule depth}\label{sec:hop}

Corollary~F of the eighth paper excludes, given a crux, every oblivious linear-time transposer whose work heads enter
every cell boundedly often~\cite[Corollary~6.1]{BoundedNesting2026}; its proof cuts the whole schedule to bounded depth
with Theorem~F. A schedule that revisits a few cells unboundedly often need not have such a cut: the deep-nesting pairs
are depth-robust if the eighth paper's Conjecture~A$'$ holds~\cite[Corollary~5.26]{BoundedNesting2026}. The eighth paper
therefore asks whether a transposer can use such a scratch to supply the information that the crux denies to the
shallow part of its schedule~\cite[Problem~10.5]{BoundedNesting2026}.

This section answers that question for codes that merge deep data only boundedly often. Theorem~F still cuts the hops
at boundedly entered cells; the hops at the other cells are charged instead of cut. Along every lineage of the code, a
chain of steps each of which continues one earlier value, one coordinate is revealed every $\tau$ deep hops, at a total
rank of $O(L/\tau)$. If the code merges deep lineages at most $s$ times along every path, the transposition matrix then
splits into a part of small rank and a part supported on balls whose radius does not depend on the schedule
(Theorems~HC and~HC$'$). So the linear crux at that radius excludes such transposers on every schedule, whatever the
number of revisited cells (Corollary~HC), and an affirmative answer to Question~9.1 of the fifth paper excludes them for
all codes (Corollary~HC$^*$). What remains concerns the code, not the schedule: routing through deep hops is limited by
their number, the number of deep cells adds nothing, and a cheap depth reduction of the code's merge structure would
suffice (\cref{prop:hop-routing,prop:capacity,prop:merge-reduction}, \cref{conj:HV}). The cold genie, the covering lemma, \cref{lem:short-paths,lem:merge-formation}, Theorem~HC,
\cref{prop:capacity} and Theorem~HC$'$ and \cref{prop:merge-reduction} for refinements of the dependency DAG are
formalized in Lean, and so are \cref{lem:expansion} for function DAGs, Corollary~HC at one size, Corollary~HC$^*$ at
one size given \cref{lem:balls}(b), and \cref{prop:hop-routing} given the count of \cref{lem:balls}(a), as their status
lines say; \cref{lem:balls} is not formalized yet.

\subsection{Deep cells and the cold genie}\label{sec:hop-cold}

Throughout this section $M$ is a machine with $k$ work tapes that is oblivious at a length $L=d^2b(d)$ with halting
time $T\le C(L+1)$, where $C\ge1$, and $\beta\ge\max(2,\beta(M))$ is an integer. We use the dependency DAG~$H_M$, the
work-head walks $x_1,\dots,x_k$ and the input-head walk~$x_0$ as in the eighth paper~\cite[Sections~2.2
and~2.4]{BoundedNesting2026}; output bit~$\ell$ is appended at step~$t_\ell$, and $\mathrm{out}(t)$ is the set of
out-neighbours of~$t$ in~$H_M$. A path $t_0\to\dots\to t_n$ is \emph{$\Gamma$-free} if $t_1,\dots,t_n\notin\Gamma$.

Fix an integer $V\ge1$. As in~\cite[Corollary~3.8]{BoundedNesting2026}, let $\mathcal R$ be the set of pairs $(h,c)$
such that the walk $x_h$ enters the cell~$c$ more than $V$ times on $[0,T)$, the \emph{deep cells}. A \emph{deep hop}
is a long hop of a head~$h$ at a cell~$c$ with $(h,c)\in\mathcal R$, and $N_{\mathcal R}$ is the number of deep hops,
counted per head. An edge of~$H_M$ is determined by its endpoints~\cite[Section~2.4]{BoundedNesting2026}; we call it
\emph{deep} if its endpoints form a deep hop of at least one head. So an edge of length at least two that is not deep is
a long hop, for each head that makes it, at a pair outside~$\mathcal R$. A long hop $(p,u)$ of a head~$h$ ends at an
entry~$u$ of its cell, since $x_h(u-1)\ne x_h(u)$, and distinct hops of~$h$ end at distinct times
(\lean{MergeDepth.card_allHops_le}); hence
\begin{equation}\label{eq:hop-NR}
  \#\{\text{deep edges of }H_M\}\;\le\;N_{\mathcal R}\;\le\;kT\;\le\;kC(L+1).
\end{equation}

\begin{lemma}[The cold genie]\label{lem:cold}
\formal{MergeDepth/HopMachine.lean}{MergeDepth.cold_genie_exists}
Let $0<\eps\le1$, $D_1=\DF(k,2V-1,\eps/(4\beta C))$, $P=\lceil4\beta C/\eps\rceil$ and $D=D_1P$. There is a set
$\Gamma_{\mathrm F}\subseteq[0,T)$ with $\beta\card{\Gamma_{\mathrm F}}\le\eps L$ such that every $\Gamma_{\mathrm F}$-free
path of~$H_M$ that contains no deep edge has at most $D-1$ edges.
\end{lemma}

\begin{outline}
Let $H'$ be the set of long hops of the work heads at pairs outside~$\mathcal R$. Each of them lies at a cell that its
head enters at most $V$ times, so $H'$ has owner-load at most $2V-1$~\cite[Corollary~3.2]{BoundedNesting2026}, and
Theorem~F~\cite[Theorem~3.6]{BoundedNesting2026}, both Lean-verified there, applied at
$\eps'=\eps/(4\beta C)\in(0,1]$, gives a cut~$F$ of $G(H')$ with at most $\eps'T$ edges and depth less than~$D_1$. Let
$\Gamma_0$ be the set of steps at which an edge of~$F$ ends, and
$\Gamma_{\mathrm F}=\Gamma_0\cup\{t<T:t\equiv P-1\pmod P\}$. Then
$\beta\card{\Gamma_{\mathrm F}}\le\beta\eps'T+\beta T/P\le\eps T/(2C)\le\eps(L+1)/2\le\eps L$.

Let $\pi=(t_0\to\dots\to t_n)$ be a $\Gamma_{\mathrm F}$-free path of~$H_M$ without deep edges. An edge of~$\pi$ of
length at least two is a long hop of some head, and since the edge is not deep, that hop lies in~$H'$. So every edge
of~$\pi$ is an edge of $G(H')$, and none lies in~$F$, because it ends outside~$\Gamma_0$. Hence $\pi$ has at most
$D_1-1$ edges of length at least two. They split~$\pi$ into at most $D_1$ runs of edges of length one; the vertices of a
run after its first are consecutive integers, none congruent to $P-1$ modulo~$P$, so a run has at most $P-1$ edges.
Thus $n\le D_1-1+D_1(P-1)=D-1$.

Formally, the work heads are walks $\N\to\Z$ and $H_M$ is replaced by the DAG of their dependency steps, the state
edges below~$T$ and the long hops of all heads (\lean{MergeDepth.DepStep}), which contains every edge of~$H_M$; a step
is deep if some head makes it as a hop at a pair of~$\mathcal R$ (\lean{MergeDepth.DeepStep}). The cut is the eighth
paper's \lean{BoundedNesting.theoremF_owner} on its \lean{BoundedNesting.lightHops}, and \lean{MergeDepth.cold_genie}
proves the path bound for every~$\eps'>0$ and~$P\ge1$, with $P\card{\Gamma_{\mathrm F}}\le P\eps'T+T$.
\end{outline}

We call such a set a \emph{cold genie at $(V,\eps)$}, and $D$ its \emph{depth}
(\lean{MergeDepth.IsColdGenie}).

This is the conversion of~\cite[Lemma~2.6]{BoundedNesting2026}, with the long-edge
conversion~\cite[Lemma~4.21]{FourthPaper2026}, restricted to paths without deep edges. It holds for every schedule,
because Theorem~F is applied only to the hops at boundedly entered cells. If every work-tape cell is entered at most $V$
times, there are no deep edges and $\Gamma_{\mathrm F}$ is the genie of Corollary~F.

\subsection{Refinements of the dependency DAG}\label{sec:hop-refine}

The genies below reveal single coordinates or single fields of a step. We state the combinatorics once, for DAGs that
refine~$H_M$.

\begin{definition}[Refinements]\label{def:refinement}
A \emph{refinement} of~$H_M$ of \emph{width}~$m$ is a finite DAG~$G$ with a map $\mathrm{st}\colon V(G)\to[0,T)$ such
that at most $m$ nodes lie at each step and every edge $u\to u'$ of~$G$ \emph{projects} to an edge
$\mathrm{st}(u)\to\mathrm{st}(u')$ of~$H_M$, together with a set of \emph{input-marked} nodes~$u$, all with
$x_0(\mathrm{st}(u))\in[0,L)$, and, for every output position~$\ell$, a set of \emph{output nodes} at step~$t_\ell$. An
edge of~$G$ is \emph{deep} if its projection is deep. For $\Gamma\subseteq[0,T)$ and a set~$Q$ of nodes, a node is
\emph{revealed} if it lies at a step of~$\Gamma$ or in~$Q$; a path $u_0\to\dots\to u_n$ of~$G$ is
\emph{$(\Gamma,Q)$-free} if $u_1,\dots,u_n$ are not revealed, and \emph{$\Gamma$-free} if it is
$(\Gamma,\emptyset)$-free. The set $I_{\Gamma,Q}(u)$ consists of the positions $x_0(\mathrm{st}(z))$ of the input-marked
nodes~$z$ that end $(\Gamma,Q)$-free paths starting at~$u$.
\end{definition}

A path of~$G$ projects to a path of~$H_M$ with the same number of edges, and a $(\Gamma,Q)$-free path projects to a
$\Gamma$-free one. Hence, if every $(\Gamma,Q)$-free path from an output node~$u$ of~$\ell$ that ends at an input-marked
node has at most $r$ edges, then $I_{\Gamma,Q}(u)\subseteq B^r(\ell)$, the set of positions read at steps reachable from
$t_\ell$ by paths of at most $r$ edges~\cite[Definition~7.2]{BoundedNesting2026}.

\paragraph{The coordinate DAG.} Suppose that $M$ has an $\F$-linear data path at length~$L$ with a representation of
width~$m$~\cite[Definition~7.1]{BoundedNesting2026}: writing $v_t\in\F^m$ for the encoded value of step~$t$ on input
$w\in\{0,1\}^L$, we have $v_t=c_t+\sum_{t'\in\mathrm{out}(t)}A_{t,t'}v_{t'}+a_tw_{x_0(t)}$, where the last term is
present only if $x_0(t)\in[0,L)$, and output bit~$\ell$ is $\varphi_\ell(v_{t_\ell})$ for an affine functional
$\varphi_\ell$ with linear part $\varphi_\ell^{\mathrm{lin}}$. The \emph{coordinate DAG} has the nodes $(t,i)$ with
$t<T$ and $i\in[m]$, where $\mathrm{st}(t,i)=t$; an edge $(t,i)\to(t',i')$ whenever $t'\in\mathrm{out}(t)$ and
$(A_{t,t'})_{ii'}\ne0$; input marks at the nodes $(t,i)$ with $x_0(t)\in[0,L)$ and $(a_t)_i\ne0$; and the output nodes
$(t_\ell,i)$ with $(\varphi^{\mathrm{lin}}_\ell)_i\ne0$. It is a refinement of width~$m$.

\paragraph{The field DAG.} The \emph{fields} of step~$t$ are the state $\sigma_t$ after the step, the symbol
$\kappa_{t,j}$ written on work tape~$j$ for $1\le j\le k$, and the output field~$o_t$, the appended bit or its absence;
each takes at most $2^\beta$ values. The \emph{predecessors} of a field of step~$t$ are $\sigma_{t-1}$ if $t\ge1$, and
$\kappa_{\prev_j(t),j}$ for every~$j$ for which $\prev_j(t)$ is defined. By the local recursion of the fourth
paper~\cite[Lemma~4.2 and its proof]{FourthPaper2026}, every field of step~$t$ is, on $\{0,1\}^L$, a function of its
predecessors and, if $x_0(t)\in[0,L)$, of the input bit $w_{x_0(t)}$. A \emph{pruning} chooses for every field~$f$ a
set $P(f)$ of its predecessors and whether $f$ \emph{keeps} its input bit, such that $f$ is still a function, on
$\{0,1\}^L$, of the fields in $P(f)$ and of the kept input bit; keeping everything is a pruning. The \emph{field DAG} of
a pruning has the edges $f\to g$ for $g\in P(f)$, input marks at the fields that keep their input bit, and the output
node $o_{t_\ell}$ for~$\ell$. It is a refinement of width $k+2$. A deep edge of it leads from a field of step~$t$ to a
field of a step $t'\le t-2$, so its head is a symbol field $\kappa_{t',j}$.

\begin{lemma}[Expansion]\label{lem:expansion}
\partformal{MergeDepth/Charging.lean}{MergeDepth.rigidity_core,MergeDepth.rigidity_finrank,MergeDepth.finrank_revealed_union_le,MergeDepth.finrank_revealed_le}{in
part~(a), the support of row~$\ell$ when the $(\Gamma,Q)$-free paths from the output nodes of~$\ell$ have no common
length bound (the formal statement assumes a bound~$D$ on their number of edges, as in all uses below); in part~(b),
that the field DAG of a pruning is a function DAG, which is the fourth paper's local recursion~\cite[Lemma~4.2]{FourthPaper2026};
for function DAGs, part~(b) is \lean{MergeDepth.FunDag.cone} with \lean{MergeDepth.FunDag.isIndexCode}}
Let $\Gamma\subseteq[0,T)$, let $Q$ be a set of nodes of the coordinate DAG or of a field DAG, and let $X$ be a set of
output positions. Let $P_{\tpos}$ be the $L\times L$ matrix with rows $e_{\tpos(\ell)}$, as
in~\cite[Section~7.2]{BoundedNesting2026}.
\begin{enumerate}[label=(\alph*)]
\item \emph{Coordinates.} For $w\in\F^L$ define $V_t(w)\in\F^m$ by the recursion of the data path, and for a node
$u=(t,i)$ let $V_u^{\mathrm{lin}}$ be the linear part of $w\mapsto V_t(w)_i$, a linear form on~$\F^L$. Let $E$ be the
span of the forms $V^{\mathrm{lin}}_u$ of the revealed nodes, and let $e_j$ also denote the form $w\mapsto w_j$. Then
$V^{\mathrm{lin}}_u\in E+\mathrm{span}\{e_j:j\in I_{\Gamma,Q}(u)\}$ for every node~$u$, and
$\dim E\le\beta\card\Gamma+\card Q$. If $M$ transposes all $d\times d$ matrices, then $P_{\tpos}=R+S$ with
$\rank R\le\beta\card\Gamma+\card Q+\card X$, where row~$\ell$ of~$S$ is zero if $\ell\in X$ and is otherwise supported
in the union $I(\ell)$ of the sets $I_{\Gamma,Q}(u)$ over the unrevealed output nodes~$u$ of~$\ell$.
\item \emph{Fields.} If $w,w'\in\{0,1\}^L$ agree on every revealed field and on $I_{\Gamma,Q}(f)$, then
$f(w)=f(w')$. If $M$ transposes all $d\times d$ matrices, then
$w\mapsto\bigl((v_t(w))_{t\in\Gamma},(f(w))_{f\in Q},(w_{\tpos(\ell)})_{\ell\in X}\bigr)$ is an index
code~\cite[Definition~4.8]{FourthPaper2026} with at most $\beta\card\Gamma+\beta\card Q+\card X$ bits for the system
$(I(\ell))_\ell$, where $I(\ell)=I_{\Gamma,Q}(o_{t_\ell})$ if $\ell\notin X$ and $o_{t_\ell}$ is not revealed, and
$I(\ell)=\emptyset$ otherwise.
\end{enumerate}
\end{lemma}

\begin{outline}
(a) The linear parts satisfy
$V^{\mathrm{lin}}_{(t,i)}=\sum_{t'\in\mathrm{out}(t)}\sum_{i'}(A_{t,t'})_{ii'}V^{\mathrm{lin}}_{(t',i')}+(a_t)_ie_{x_0(t)}$,
the last term present only if $x_0(t)\in[0,L)$; the nonzero terms belong to the out-edges of $(t,i)$ and to its input
mark. We prove the first claim by induction on~$t$. If $u=(t,i)$ is input-marked, the path of length zero shows
$x_0(t)\in I_{\Gamma,Q}(u)$. For an edge $u\to z$, either $z$ is revealed and $V^{\mathrm{lin}}_z\in E$, or a
$(\Gamma,Q)$-free path from~$z$, preceded by $u\to z$, is $(\Gamma,Q)$-free, so
$I_{\Gamma,Q}(z)\subseteq I_{\Gamma,Q}(u)$ and the induction hypothesis applies to~$z$.

The forms of the nodes of~$Q$ span at most $\card Q$ dimensions. On binary inputs the formal values agree with the
encoded values, $V_t(w)=v_t$, by induction on~$t$, so the linear map $\Phi\colon w\mapsto(V_t(w)-V_t(0))_{t\in\Gamma}$
takes at most $2^{\beta\card\Gamma}$ values on $\{0,1\}^L$. A linear map of rank~$\rho$ takes at least $2^\rho$ values
there: it is injective on the binary vectors supported on $\rho$ coordinates whose columns are independent, because the
difference of two such vectors has coefficients in $\{-1,0,1\}$, not all zero. The forms of the nodes at steps
of~$\Gamma$ span the row space of~$\Phi$, so $\dim E\le\beta\card\Gamma+\card Q$.

If $M$ transposes, $Y_\ell(w)=\varphi_\ell(V_{t_\ell}(w))$ is an affine function on~$\F^L$ that equals $w_{\tpos(\ell)}$
on $\{0,1\}^L$. Evaluating it at~$0$ and at the unit vectors shows that its linear part
$\sum_i(\varphi^{\mathrm{lin}}_\ell)_iV^{\mathrm{lin}}_{(t_\ell,i)}$ is the form $e_{\tpos(\ell)}$. A term with
$(\varphi^{\mathrm{lin}}_\ell)_i\ne0$ belongs to an output node; it lies in~$E$ if the node is revealed, and in
$E+\mathrm{span}\{e_j:j\in I(\ell)\}$ otherwise. So $e_{\tpos(\ell)}=r_\ell+s_\ell$ with $r_\ell\in E$ and $s_\ell$
supported in $I(\ell)$; for $\ell\in X$ put $r_\ell=e_{\tpos(\ell)}$ and $s_\ell=0$. The matrices with rows $r_\ell$
and~$s_\ell$ are as claimed, since the rows $r_\ell$ lie in $E+\mathrm{span}\{e_{\tpos(\ell)}:\ell\in X\}$.

(b) The first claim is the cone lemma of the fourth paper~\cite[Lemma~4.4 and Remark~4.6]{FourthPaper2026}, with the same
proof by induction on the step of~$f$: $f$ is a function of the fields in $P(f)$ and of its kept input bit, whose
position lies in $I_{\Gamma,Q}(f)$; a revealed field in $P(f)$ takes the same value on $w$ and~$w'$; and an unrevealed
$g\in P(f)$ satisfies $I_{\Gamma,Q}(g)\subseteq I_{\Gamma,Q}(f)$, as in~(a). Now let $w,w'$ have the same image under the
map and agree on $I(\ell)$. They agree on every revealed field, since $v_t$ determines the fields of step~$t$. If
$\ell\in X$, the bit $w_{\tpos(\ell)}$ is part of the image. Otherwise $o_{t_\ell}$ takes the same value on $w$ and~$w'$,
directly or by the first claim, and its bit is $w_{\tpos(\ell)}$. A step value and a field take at most $2^\beta$ values
each.

Formally, the decomposition is \lean{MergeDepth.rigidity_core} and \lean{MergeDepth.rigidity_finrank}, for a linear
network whose output~$\ell$ is an affine function of the values at a set of output nodes; the bound on the revealed rank
is \lean{MergeDepth.finrank_revealed_union_le} with \lean{MergeDepth.finrank_revealed_le} for single coordinates, and
\lean{MergeDepth.finrank_revealed_le_of_values} with \lean{MergeDepth.finrank_revealed_coord_le} for whole steps.
\end{outline}

For whole steps, the rank bound of~(a) is the counting of the eighth paper's per-machine rigidity
lemma~\cite[Lemma~7.3]{BoundedNesting2026}. Its Lean statement charges $\beta$ per revealed node of a network, so applied to
the coordinate DAG it would give only $\beta(m\card\Gamma+\card Q)$; the bound for mixed genies, $\beta$ per revealed step
and one per revealed coordinate, is the elementary extension proved above.

\subsection{Hop charging along transport trees}\label{sec:hop-HC}

\begin{definition}[Transport trees and deep formations]\label{def:transport}
\formal{MergeDepth/Formation.lean}{MergeDepth.IsCopy,MergeDepth.copyForest,MergeDepth.root,MergeDepth.DeepFormation}
Let $G$ be a refinement. A \emph{copy} is a node that has exactly one out-edge and is not input-marked; every other node
is a \emph{formation}. Following the out-edges of copies from a node~$u$ leads, as steps decrease, to a formation
$\mathrm{root}(u)$, with $\mathrm{root}(u)=u$ for a formation; the edges followed form the \emph{chain} of~$u$. The
\emph{transport tree} of a formation~$p$ is the set of nodes~$u$ with $\mathrm{root}(u)=p$, and its edges are the
out-edges of its copies. Transport trees partition $V(G)$, and the head of an edge of a transport tree lies in the same
tree. A transport tree is \emph{$\tau$-tall} if the chain of one of its nodes contains at least $\tau$ deep edges. A
formation is \emph{deep} if it is the tail of a deep edge or its transport tree contains a deep edge. For
$\Gamma_{\mathrm F}\subseteq[0,T)$, the \emph{formation depth} at an output position~$\ell$ is the largest number of
deep formations on a $\Gamma_{\mathrm F}$-free path of~$G$ from an output node of~$\ell$.
\end{definition}

In the coordinate DAG, a copy $u$ with out-neighbour~$u'$ satisfies $V^{\mathrm{lin}}_u=\lambda V^{\mathrm{lin}}_{u'}$
for some $\lambda\ne0$. Storing a value in a cell, carrying it in the state, moving it between cells and multicasting it
create no formations; combining it with anything else does.

\begin{theorem}[Theorem HC: hop charging]\label{thm:HC}
\formal{MergeDepth/HopMachine.lean}{MergeDepth.hop_charging_HC}
Let $M$ have an $\F$-linear data path at length~$L$ with a representation of width~$m$, let $\Gamma_{\mathrm F}$ be a
cold genie at $(V,\eps)$ with depth~$D$, let $\tau\ge1$ and $s\ge0$ be integers, and let $X$ be the set of output
positions at which the coordinate DAG has formation depth greater than~$s$. If $M$ transposes all $d\times d$ matrices,
then
\[
  P_{\tpos}=R+S,\qquad\rank R\le\eps L+\frac{mkC(L+1)}{\tau}+\card X,
\]
and row~$\ell$ of~$S$ is zero for $\ell\in X$ and supported in $B^{(s+1)\tau D}(\ell)$ otherwise.
\end{theorem}

\begin{outline}
Give every copy its out-edge as lineage edge and no other node a lineage edge; this is a lineage assignment for every
$\Gamma_{\mathrm F}$ (\cref{def:lineage,rem:hop-examples}). By \cref{lem:merge-formation}, its merge depth is at most the
formation depth at every output position, so \cref{thm:HCprime}(a) applies with the same~$s$ and~$X$.
\end{outline}

\begin{remark}[A direct proof]\label{rem:HC-direct}
Theorem~HC also has a direct proof with a simpler genie, the set~$Q_0$ of roots of the $\tau$-tall transport trees, at
the slightly larger radius $r_0=((s+1)\tau+1)D$: $\card{Q_0}\le mN_{\mathcal R}/\tau$, and every
$(\Gamma_{\mathrm F},Q_0)$-free path that starts at an unrevealed output node of a position $\ell\notin X$ and ends at an
input-marked node has fewer than $r_0$ edges. This proof is not formalized.

\emph{Counting tall trees.} For every tall tree fix a node whose chain contains at least $\tau$ deep edges. The tails of
the edges of a chain lie at distinct steps, so the deep edges of a fixed chain project to distinct deep edges of~$H_M$.
Fix a deep edge $t\to t'$ of~$H_M$. An edge of a fixed chain that projects to it has its tail at step~$t$ and its head at
step~$t'$, and it lies in the tree of its tail, which contains its head. Trees are disjoint, so the fixed chains of
distinct trees that use $t\to t'$ have distinct heads at step~$t'$, at most $m$ of them. Counting the pairs (tall tree,
deep edge on its fixed chain) gives $\tau\card{Q_0}\le m\cdot\#\{\text{deep edges of }H_M\}\le mN_{\mathcal R}$
by~\eqref{eq:hop-NR}.

\emph{Paths.} Let $\pi=(u_0\to\dots\to u_n)$ be $(\Gamma_{\mathrm F},Q_0)$-free, with $u_0$ an unrevealed output node of
a position $\ell\notin X$ and $u_n$ input-marked. First, $\pi$ meets no tall tree. Indeed, a vertex $u_i$ of a tall tree
is not its root, which is revealed, while $u_0$ is unrevealed and so are the later vertices; so $u_i$ is a copy, $\pi$
continues along its out-edge inside the tree, and by the same argument through copies of that tree to its end; but
$u_n$ is input-marked, hence not a copy. Second, $\pi$ is $\Gamma_{\mathrm F}$-free, so it contains at most $s$ deep
formations, say at positions $i_1<\dots<i_{s'}$; put $i_0=0$ and $i_{s'+1}=n$. Consider the piece of~$\pi$ from
$u_{i_a}$ to $u_{i_{a+1}}$ and a deep edge $u_i\to u_{i+1}$ of it whose tail is not a deep formation; only the edge
leaving $u_{i_a}$ can fail this condition. A formation with a deep out-edge is deep, so $u_i$ is a copy, and the edge lies
on the chain of~$u_i$, in a tree that is not tall. From there $\pi$ follows that chain through copies up to the root,
which is a deep formation because its tree contains a deep edge, hence is $u_{i_{a+1}}$; or $\pi$ ends before. So all
deep edges of the piece, except possibly the one leaving $u_{i_a}$, lie on one chain of a tree that is not tall, and there
are fewer than $\tau$ of them. Each piece thus has at most $\tau$ deep edges, and $\pi$ has at most $h\le(s+1)\tau$.
Removing them leaves at most $h+1$ runs; the vertices of a run after its first are not revealed, so a run projects to a
$\Gamma_{\mathrm F}$-free path of~$H_M$ without deep edges and has at most $D-1$ edges by \cref{lem:cold}. Hence
$n\le h+(h+1)(D-1)\le r_0-1$.

\emph{Expansion.} Apply \cref{lem:expansion}(a) with $\Gamma=\Gamma_{\mathrm F}$ and $Q=Q_0$. For $\ell\notin X$ the set
$I(\ell)$ lies in $B^{r_0-1}(\ell)\subseteq B^{r_0}(\ell)$ by the bound on paths, and
$\rank R\le\beta\card{\Gamma_{\mathrm F}}+\card{Q_0}+\card X\le\eps L+mN_{\mathcal R}/\tau+\card X$; apply
\eqref{eq:hop-NR}.
\end{remark}

Theorem~HC charges every deep hop on a long chain of copies to the root of the chain: storage, transport and multicast
through deep cells cost rank $O(N_{\mathcal R}/\tau)$ on every schedule, and only formations at deep cells count. The
direct proof applies verbatim to field DAGs, with \cref{lem:expansion}(b) in place of~(a); the revealed roots are then fields of
any kind, a state or an output field among them, and cost up to $\beta$ bits each.

\subsection{Hop charging along lineages}\label{sec:hop-HCprime}

Theorem~HC counts every formation at a deep cell, so it does not cover a deep cell that accumulates fresh data,
$\sigma\leftarrow\sigma+x$, at each visit. Such an update keeps a single lineage, and lineages can be checkpointed.

\begin{definition}[Lineages and merge depth]\label{def:lineage}
Let $G$ be a refinement and $\Gamma_{\mathrm F}\subseteq[0,T)$. A \emph{lineage assignment} for $\Gamma_{\mathrm F}$
chooses, for some nodes~$u$, an out-edge $e_u$ of~$u$ in~$G$, the \emph{lineage edge} of~$u$, such that every
$\Gamma_{\mathrm F}$-free path whose first edge leaves~$u$ and differs from $e_u$ contains no deep edge. The nodes with a
lineage edge are \emph{continuations}; all other nodes, including those without out-edges, are \emph{merge nodes}.
Following lineage edges from a node leads, as steps decrease, to a merge node, its \emph{lineage root}; the nodes with
lineage root~$p$ form the \emph{lineage tree} of~$p$, a tree with root~$p$ whose edges are the lineage edges of its
continuations. Lineage trees partition $V(G)$. A node $u_i$ of a path $\pi=(u_0\to\dots\to u_n)$ of~$G$ is a \emph{deep
merge} of~$\pi$ if it is a merge node and
\begin{enumerate}
\item[($\alpha$)] $i<n$ and the edge $u_i\to u_{i+1}$ is deep, or
\item[($\beta$)] for some $j<i$ the edge $u_j\to u_{j+1}$ is deep, and for $j\le j'<i$ the node $u_{j'}$ is a
continuation and $u_{j'}\to u_{j'+1}$ is its lineage edge.
\end{enumerate}
The \emph{merge depth} at an output position~$\ell$ is the largest number of deep merges of a
$\Gamma_{\mathrm F}$-free path of~$G$ from an output node of~$\ell$. The refinement has \emph{merge depth at most~$s$},
or is \emph{$s$-round}, if its merge depth is at most~$s$ at every output position.
\end{definition}

Lineage edges are edges of~$G$, never input marks. A merge node with out-edges typically merges two deep lineages, or
deep content with data derived from deep cells. Whether a node is a deep merge depends on the path, which must leave it
along a deep edge or arrive through a deep lineage edge followed only by lineage edges. Counting every merge node whose
lineage tree contains a deep edge would not satisfy \cref{lem:merge-formation} below: such a node need not be a deep
formation, since its transport tree may contain no deep edge, and a run of cold edges can pass through many such nodes
without using their lineage trees.

\begin{remark}[Examples]\label{rem:hop-examples}
Assigning to every copy its out-edge, and to no other node a lineage edge, is a lineage assignment for every
$\Gamma_{\mathrm F}$. A deep-cell update that combines the previous content of the cell with data whose
$\Gamma_{\mathrm F}$-free paths reach no deep edge, such as fresh input bits, is a continuation. If every tail of a deep
edge is a continuation with that edge as lineage edge, and no lineage root of a tree containing a deep edge reaches a deep
edge by a $\Gamma_{\mathrm F}$-free path, then ($\alpha$) never applies, nothing after a deep merge of type~($\beta$) is
deep, and every path has at most one deep merge. This covers deep cells that only store, move, multicast or accumulate
data not derived from deep cells, and pure data movement. Relays with boundedly many mixing steps, or butterflies of
bounded depth, have bounded merge depth. A side edge into a step of $\Gamma_{\mathrm F}$ may be deep, as the one-edge path
into a revealed step is not $\Gamma_{\mathrm F}$-free; the paths that matter below never use one.
\end{remark}

\begin{lemma}[Covering]\label{lem:covering}
\formal{MergeDepth/Covering.lean}{MergeDepth.covering}
Let $\mathcal T$ be a finite forest whose edges are directed towards the roots, some of them marked, and let $\tau\ge1$.
There is a set~$Q$ of heads of marked edges, with $\tau\card Q$ at most the number of marked edges, such that every walk
$x_0\to x_1\to\dots\to x_n$ along edges of~$\mathcal T$ with $x_1,\dots,x_n\notin Q$ crosses fewer than $\tau$ marked
edges.
\end{lemma}

\begin{outline}
The \emph{marked height} of a node is the number of marked edges on its walk to its root; along an edge it drops by one
if the edge is marked and stays the same otherwise. Sort the heads of the marked edges into $\tau$ classes by their
marked heights modulo~$\tau$, and let~$Q$ be a smallest class. There are at most as many heads as marked edges, so
$\tau\card Q$ is at most the number of marked edges. A walk that crosses $\tau$ marked edges meets, after its first
node, heads of marked edges at $\tau$ consecutive marked heights, one in each class, so one of them lies in~$Q$.
\end{outline}

\begin{lemma}[Few deep merges, short paths]\label{lem:short-paths}
\formal{MergeDepth/HopFields.lean}{MergeDepth.short_paths}
Let $G$ be a refinement of width~$m$, let $\Gamma_{\mathrm F}$ be a cold genie at $(V,\eps)$ with depth~$D$, fix a
lineage assignment for~$\Gamma_{\mathrm F}$, and let $\tau\ge1$. Apply \cref{lem:covering} to the forest of lineage
edges, with its deep edges marked, and let $Q$ be the set obtained. Then $\card Q\le mN_{\mathcal R}/\tau$, and every
$(\Gamma_{\mathrm F},Q)$-free path of~$G$ with $s'$ deep merges has at most $(s'+1)\tau-1$ deep edges and fewer than
$(s'+1)\tau D$ edges.
\end{lemma}

\begin{outline}
\emph{Counting.} A deep lineage edge projects to a deep edge $t\to t'$ of~$H_M$ and has its tail at step~$t$; every node
has at most one lineage edge, and at most $m$ nodes lie at step~$t$. So there are at most $m$ times as many deep lineage
edges as deep edges of~$H_M$, at most $mN_{\mathcal R}$ by~\eqref{eq:hop-NR}, and \cref{lem:covering} gives
$\card Q\le mN_{\mathcal R}/\tau$.

\emph{Deep edges.} Let $\pi=(u_0\to\dots\to u_n)$ be $(\Gamma_{\mathrm F},Q)$-free. If $\pi$ leaves a continuation~$u_i$
by an edge other than its lineage edge, then the rest $u_i\to\dots\to u_n$ contains no deep edge, since it is
$\Gamma_{\mathrm F}$-free. Walk along~$\pi$ and keep the current \emph{stretch}: the longest final segment of the part
walked so far that follows lineage edges. Its vertices after the first are vertices of~$\pi$ after~$u_0$, hence not
in~$Q$, so by \cref{lem:covering} it contains fewer than $\tau$ deep edges. By induction on the length of the rest of the
path, the deep edges of the rest together with those of the current stretch number at most $\tau-1$ plus $\tau$ times the
deep merges of the rest, where a merge node reached through the current stretch counts as of type~($\beta$) if the stretch
contains a deep edge. Indeed, at a merge node~$u_i$ the stretch ends: if the edge $u_i\to u_{i+1}$ is deep or the stretch
contains a deep edge, then $u_i$ is a deep merge, and its share~$\tau$ pays for the fewer than $\tau$ deep edges of the
stretch and the edge leaving~$u_i$; otherwise the stretch contains no deep edge. At a continuation that $\pi$ leaves by
its lineage edge the stretch grows by that edge, and at a continuation that $\pi$ leaves by another edge no deep edge
follows. Starting with the stretch $(u_0)$, the path has at most $h\le\tau-1+\tau s'=(s'+1)\tau-1$ deep edges.

\emph{Length.} Removing the deep edges leaves at most $h+1$ runs. The vertices of a run after its first are not revealed,
so a run projects to a $\Gamma_{\mathrm F}$-free path of~$H_M$ without deep edges and has at most $D-1$ edges by the
definition of the depth. Hence $n\le h+(h+1)(D-1)=(h+1)D-1<(s'+1)\tau D$.
\end{outline}

\begin{theorem}[Theorem HC$'$: hop charging along lineages]\label{thm:HCprime}
\partformal{MergeDepth/HopMachine.lean}{MergeDepth.hop_charging_machine}{in part~(b), that the field DAG of a pruning
is a function DAG that refines the dependency DAG, which is the fourth paper's local recursion; for such refinements,
part~(b) is \lean{MergeDepth.hop_charging_fields}}
Let $G$ be either the coordinate DAG of an $\F$-linear data path of~$M$ at length~$L$ of width~$m$, or the field DAG of
a pruning, with $m=k+2$. Let $\Gamma_{\mathrm F}$ be a cold genie at $(V,\eps)$ with depth~$D$, fix a lineage
assignment for~$\Gamma_{\mathrm F}$, let $\tau\ge1$ and $s\ge0$ be integers, put $r=(s+1)\tau D$, and let $X$ be
the set of output positions at which $G$ has merge depth greater than~$s$. There is a set~$Q$ of at most
$mN_{\mathcal R}/\tau$ nodes of~$G$ such that every $(\Gamma_{\mathrm F},Q)$-free path of~$G$ from an output node of a
position $\ell\notin X$ has fewer than $r$ edges. Suppose that $M$ transposes all $d\times d$ matrices.
\begin{enumerate}[label=(\alph*)]
\item If $G$ is the coordinate DAG, then $P_{\tpos}=R+S$ with $\rank R\le\eps L+mkC(L+1)/\tau+\card X$ and row~$\ell$
of~$S$ supported in $B^r(\ell)$ for every~$\ell$.
\item If $G$ is a field DAG, there is an index code with at most $\eps L+\beta(k+2)kC(L+1)/\tau+\card X$ bits for a
system $(I(\ell))_\ell$ with $I(\ell)\subseteq B^r(\ell)$ for every~$\ell$.
\end{enumerate}
\end{theorem}

\begin{outline}
Let $Q$ be the set of \cref{lem:short-paths}. A $(\Gamma_{\mathrm F},Q)$-free path from an output node of $\ell\notin X$ is
$\Gamma_{\mathrm F}$-free, so it has at most $s$ deep merges and, by \cref{lem:short-paths}, at most $r-1$ edges. Hence
the sets $I(\ell)$ of \cref{lem:expansion}, applied with $\Gamma=\Gamma_{\mathrm F}$, $Q$ and~$X$, lie in
$B^{r-1}(\ell)\subseteq B^r(\ell)$ for $\ell\notin X$, and they are empty for $\ell\in X$ in case~(b). The rank in~(a) is
at most $\beta\card{\Gamma_{\mathrm F}}+\card Q+\card X\le\eps L+mN_{\mathcal R}/\tau+\card X$, and the code in~(b) has at
most $\beta\card{\Gamma_{\mathrm F}}+\beta\card Q+\card X\le\eps L+\beta(k+2)N_{\mathcal R}/\tau+\card X$ bits;
apply~\eqref{eq:hop-NR}. Formally, part~(a) is proved for a linear network of width~$m$ whose edges are dependency steps
of the work-head walks, whose steps take at most $2^\beta$ values on binary inputs, and a cold genie in the sense of
\lean{MergeDepth.IsColdGenie}. \lean{MergeDepth.hop_charging_steps} proves it with the rank bound
$\beta\card{\Gamma_{\mathrm F}}+q+\card X$, where $\tau q$ is at most $m$ times the number of deep edges; the
coordinate DAG is the network \lean{MergeDepth.coordNet}, whose values are the coordinates of the step values
(\lean{MergeDepth.coord_val}) and whose paths and balls project to those of~$H_M$
(\lean{MergeDepth.coord_paths}, \lean{MergeDepth.coord_reach_subset}); \lean{MergeDepth.hop_charging} combines
\lean{MergeDepth.paths_short} with the decomposition \lean{MergeDepth.rigidity_finrank} of \cref{lem:expansion}(a).
\end{outline}

Theorem~HC$'$ reduces the depth of the code, not of the schedule: every output has bounded depth in~$G$ once
$\Gamma_{\mathrm F}$ and~$Q$ are revealed, as for the coordinate genies of~\cite[Remark~7.6(c)]{BoundedNesting2026},
while paths of~$H_M$, which need not follow the dependencies of the code, may stay long. In~(a) the bound is on the rank
and does not depend on~$\F$; counted in bits, a revealed coordinate would cost up to
$\min(\lceil\log_2\card\F\rceil,\beta)$.

\begin{lemma}[Merge depth and formation depth]\label{lem:merge-formation}
\formal{MergeDepth/Formation.lean}{MergeDepth.dm_le_countP_deepFormation}
Let $G$ be a refinement and $\Gamma_{\mathrm F}\subseteq[0,T)$, and fix a lineage assignment for $\Gamma_{\mathrm F}$ in
which every copy is a continuation. Then every path of~$G$ contains at least as many deep formations as it has deep
merges. In particular, the merge depth at every output position is at most the formation depth there.
\end{lemma}

\begin{outline}
A copy has only one out-edge, which is therefore its lineage edge, and a merge node is not a copy, hence a formation.
Walk along a path with a credit for the current stretch of lineage edges. When the path crosses a deep lineage edge out
of a formation, that formation is a deep formation of the path and credits the stretch. When it crosses one out of a
copy, the next formation on the path lies on the chain of that copy, so it is the root of a transport tree containing a
deep edge, a deep formation; until it is reached, the deep copy edge lies in the transport tree of the current node. A
deep merge of type~($\alpha$) is a formation with a deep out-edge, hence a deep formation itself; one of type~($\beta$)
ends a stretch that contains a deep lineage edge, and is paid by the credit of that stretch or, if the deep edge left a
copy and no formation has followed, is itself the root of the transport tree that contains it. A side edge ends the
stretch, and stretches are disjoint, so no deep formation pays twice.
\end{outline}

So Theorem~HC$'$ applies, with the same~$s$, to every code to which Theorem~HC applies, at the smaller radius
$(s+1)\tau D$ instead of $((s+1)\tau+1)D$, and it covers accumulators of fresh data as well. It is the form used below.

\subsection{Exclusion of bounded merge depth}\label{sec:hop-excl}

\begin{corollary}[Corollary HC: the linear crux excludes bounded merge depth]\label{cor:HC}
\partformal{MergeDepth/HopMachine.lean}{MergeDepth.cor_HC_step,MergeDepth.cor_HC_step_formation}{the quantification over
machines and sizes, and the linear crux as a statement about the class $\Sball_{k,C,\beta,r}$; the Lean theorems are the
argument at one size: the ball system of radius~$r$ has minrank less than $(\rho-\eta')L$, $\eta'=(\rho-\eps)/2$, for
merge depth and for formation depth}
Let $k,m,V\ge1$, $\beta\ge2$ and $s\ge0$ be integers, $C\ge1$, $\F$ a field, $0<\eps<\rho$ with $\eps\le1$, and
$0\le\eta<(\rho-\eps)/8$. Put
\[
  \tau=\Bigl\lceil\frac{8mkC}{\rho-\eps}\Bigr\rceil,\qquad
  D=\DF\bigl(k,2V-1,\eps/(4\beta C)\bigr)\cdot\lceil4\beta C/\eps\rceil,\qquad r=(s+1)\tau D,
\]
and suppose that $\LinCrux_\F(\Sball_{k,C,\beta,r},\rho)$ holds~\cite[Definition~7.4]{BoundedNesting2026}. Then for every
machine~$M$ with $k$ work tapes and value width at most~$\beta$ there is~$d_0$ such that for no $d\ge d_0$ do the
following hold together at $L=d^2b(d)$:
\begin{enumerate}[label=(\roman*)]
\item $M$ is oblivious at length~$L$ with halting time at most $C(L+1)$ and transposes all $d\times d$ matrices;
\item $M$ has an $\F$-linear data path at length~$L$ with a representation of width at most~$m$ whose coordinate DAG,
for some cold genie at $(V,\eps)$ and some lineage assignment for it, has merge depth at most~$s$ at all but at most
$\eta L$ output positions.
\end{enumerate}
The same holds with merge depth replaced by formation depth.
\end{corollary}

\begin{outline}
Let $d_0$ be given by the linear crux for $\eta'=(\rho-\eps)/2$, and suppose that (i) and~(ii) hold at some $d\ge d_0$,
with a representation of width $m'\le m$, a cold genie $\Gamma_{\mathrm F}$ and a lineage assignment; let $X$ be the set
of output positions at which the merge depth exceeds~$s$, so $\card X\le\eta L$. \Cref{thm:HCprime}(a), with $m'$ in place
of~$m$, gives $P_{\tpos}=R+S$ with row~$\ell$ of~$S$ supported in $B^r(\ell)$ and, as $\tau\ge8mkC/(\rho-\eps)$ and
$L\ge1$,
\[
  \rank R\le\eps L+\frac{(\rho-\eps)(L+1)}8+\eta L<\eps L+\frac{(\rho-\eps)(2L+1)}8\le\eps L+\frac{(\rho-\eps)L}2
  =(\rho-\eta')L.
\]
The ball system $(B^r(\ell))_\ell$ of~$M$ belongs to $\Sball_{k,C,\beta,r}$, since $M$ is oblivious with $k$ work tapes,
value width at most~$\beta$ and halting time at most $C(L+1)$, and appends exactly $L$ bits. Its minrank over~$\F$ is at
most $\rank R<(\rho-\eta')L$, contrary to the linear crux. For formation depth, take the lineage assignment that gives every copy its out-edge
(\cref{rem:hop-examples}); by \cref{lem:merge-formation} its merge depth is at most the formation depth.
\end{outline}

Given the linear crux at one fixed radius, an oblivious linear-time transposer with a linear code must therefore, for
every~$s$ and on every schedule, merge deep lineages more than $s$ times along some path from each of more than
$(\rho-\eps)L/8$ outputs.

\begin{remark}[Parameters and special cases]\label{rem:HC-scope}
(a) The depth in Corollary~F depends only on $k$, $V$, $\beta$, $C$ and~$\eps$; the radius~$r$ here depends also on
$\rho-\eps$, $m$ and~$s$, through~$\tau$. Excluding machines of every width, or of every merge depth, needs the linear crux
along a sequence of radii. (b) If every work-tape cell is entered at most $V$ times, there are no deep edges, no deep
merges and $Q=\emptyset$, and row~$\ell$ of~$S$ is supported in the cone of the whole-step genie~$\Gamma_{\mathrm F}$, of
full depth less than~$D$. The argument then needs only a minrank bound for the cone systems of the class
$\mathcal S_{k,C,\beta,D}$~\cite[Definition~4.8]{FourthPaper2026}: it is Corollary~F restricted to linear codes. Stated
with balls, Corollary~HC contains only the ball form of Corollary~F, whose hypothesis
$\LinCrux_\F(\Sball_{k,C,\beta,D},\rho)$ is stronger, as ball systems dominate cone
systems~\cite[Remark~7.6(a)]{BoundedNesting2026}. (c) No depth reduction of the deep part is used, so the corollary applies
to the deep-nesting pairs as to every other schedule.
\end{remark}

For all codes the crux has to be applied to relations that are not cone relations of whole-step genies. The fifth
paper's structural theorem extends to them, because its proof uses the genie only through short paths.

\begin{lemma}[Relations inside balls]\label{lem:balls}
\notformalized
Let $D'\ge0$ be an integer and $K=\sum_{i=0}^{D'}(k+1)^i$. Suppose that $d\ge2$ and that $M$ appends exactly $L$ output
bits, and let $N$ be a recovery relation on $[0,L)$ with $N(q)\subseteq B^{D'}(\tpos(q))$ for every~$q$.
\begin{enumerate}[label=(\alph*)]
\item The conclusions of~\cite[Theorem~3.3, Corollary~3.7 and Corollary~3.8]{DirectedRecovery2026} hold for~$N$, with
$D'$ in place of the full depth of the genie and with the same constants. In particular, at most
$(4C(K+1)(k+2)(k+6)+3)L/\log_2d$ positions~$p$ satisfy $p\in N(p)$.
\item If Question~9.1 of~\cite{DirectedRecovery2026} has an affirmative answer for every fixed degree bound, then for
every $\delta>0$ there is~$d_1$, depending only on $k$, $D'$, $C$ and~$\delta$, such that for $d\ge d_1$ every storage
code~$\mathcal C$ for~$N$~\cite[Definition~2.1]{DirectedRecovery2026} over a finite alphabet~$\mathcal A$ with
$\card{\mathcal A}\ge2$ satisfies $\log_{\card{\mathcal A}}\card{\mathcal C}\le(\frac12+\delta)L$.
\end{enumerate}
\end{lemma}

\begin{informalproof}
(a) In the fifth paper the relation is $N(q)=I_\Gamma(\Pi(q))$ for a genie~$\Gamma$ of full depth at most~$D$, where
$\Pi=\tpos$ is the transpose permutation of~\cite[Section~3.1]{DirectedRecovery2026}, and $B=\sum_{i\le D}(k+1)^i$. The
proofs there of Lemma~3.6, Theorem~3.3 and Corollaries~3.7 and~3.8 use the relation only through Lemma~3.5 and, in
Corollary~3.8, through the bound $\card{N(q)}\le B$; Lemma~3.4 involves no relation. The proof of Lemma~3.5 uses the
genie in one place only: for every edge $L_pR_q$ of a matching it chooses a $\Gamma$-free path of~$H_M$ with at most $D$
edges from the step that appends output position $\Pi(q)$ to a step that reads~$p$. Afterwards it uses only that the
chosen paths are paths of~$H_M$ with at most $D$ edges, and that every vertex of~$H_M$ has in-degree at most $k+1$, so
that a vertex lies on the chosen paths of at most $B$ output steps. For our~$N$ a path with at most $D'$ edges from
$t_{\tpos(q)}$ to a step that reads~$p$ exists for every $p\in N(q)$, by the definition of $B^{D'}(\tpos(q))$; and
$\card{N(q)}\le K$, because at most $K$ steps are reachable from a step by paths of at most $D'$ edges, out-degrees being
at most $k+1$, and each step reads at most one position. So the proofs apply verbatim, with $D'$ and~$K$ in place of $D$
and~$B$. The last sentence of~(a) is Corollary~3.7, with the constant~$A_1$ of Lemma~3.6.

(b) This is the proof of~\cite[Proposition~9.2]{DirectedRecovery2026}, with~(a) in place of the statements for cone
relations. Let $\delta>0$, put $M_\delta=\lceil4K/\delta\rceil$ and $\Delta=\max(K,M_\delta)$, and suppose that the claim
fails: there are machines as in the statement at sizes $d_j\to\infty$, relations~$N_j$ and codes~$\mathcal C_j$ over
alphabets~$\mathcal A_j$ with $\log_{\card{\mathcal A_j}}\card{\mathcal C_j}>(\frac12+\delta)L_j$. Let $H_j$ be the set of
coordinates~$p$ whose left vertex $L_p$ has degree greater than $M_\delta$ in the incidence
graph~\cite[Definition~2.3]{DirectedRecovery2026}. The left degrees sum to $\sum_q\card{N_j(q)}\le KL_j$, so
$\card{H_j}\le\delta L_j/4$. Let $Q_j$ be the complement of~$H_j$ and $N'_j$ the relation induced on~$Q_j$. Its incidence
graph is the induced graph $E_{Q_j}$, of maximum degree at most~$\Delta$, and by~(a), Theorem~3.3 applied with $Q=Q_j$,
for every fixed~$r$ at most $O(L_j/\log d_j)=o(\card{Q_j})$ coordinates of~$Q_j$ have their two ports at distance at
most~$r$; so $(N'_j)$ satisfies the paired-distance condition. Fixing the coordinates in~$H_j$ and restricting to~$Q_j$
gives storage codes for~$N'_j$~\cite[Lemma~2.2]{DirectedRecovery2026}, one of which, $\mathcal C'_j$, has
$\log_{\card{\mathcal A_j}}\card{\mathcal C'_j}\ge\log_{\card{\mathcal A_j}}\card{\mathcal C_j}-\card{H_j}
>(\frac12+\frac{3\delta}4)L_j\ge(\frac12+\frac{3\delta}4)\card{Q_j}$. This contradicts the affirmative answer to
Question~9.1 for the degree bound~$\Delta$.
\end{informalproof}

\begin{corollary}[Corollary HC$^*$: Question~9.1 excludes bounded merge depth for all codes]\label{cor:HCstar}
\partformal{MergeDepth/HCStar.lean}{MergeDepth.cor_HCstar_step}{\cref{lem:balls}(b), which the formal statement assumes at
the size in question; that the field DAG of a pruning is a function DAG that refines the dependency DAG, as in
\cref{thm:HCprime}(b); and the last sentence}
Suppose that Question~9.1 of~\cite{DirectedRecovery2026} has an affirmative answer for every fixed degree bound; the
binary alphabet suffices. Let $k,V\ge1$ and $s\ge0$ be integers, $C\ge1$, and let $0<\eps\le1$ and $\eta\ge0$ satisfy
$\eps+\eta<\frac12$. Then for every machine~$M$ with $k$ work tapes there is~$d_0$ such that for no $d\ge d_0$ do the
following hold together at $L=d^2b(d)$:
\begin{enumerate}[label=(\roman*)]
\item $M$ is oblivious at length~$L$ with halting time at most $C(L+1)$ and transposes all $d\times d$ matrices;
\item for some pruning, the field DAG of~$M$ at length~$L$ has, for some cold genie at $(V,\eps)$, taken with
$\beta=\max(2,\beta(M))$, and some lineage assignment for it, merge depth at most~$s$ at all but at most $\eta L$ output
positions.
\end{enumerate}
In particular, for all $V\ge1$ and $s\ge0$ there is~$d_0$ such that at no $d\ge d_0$ at which (i) holds does a pruning
with a lineage assignment for the empty genie give the field DAG merge depth at most~$s$ with respect to the empty genie.
\end{corollary}

\begin{outline}
Put $\beta=\max(2,\beta(M))$, $\gamma=\frac12-\eps-\eta>0$ and $\tau=\lceil4\beta(k+2)kC/\gamma\rceil$, let $D$ be the
depth of the cold genies at $(V,\eps)$, and put $r=(s+1)\tau D$; these depend only on $M$, $V$, $s$, $C$, $\eps$
and~$\eta$. Let $d_1$ be given by \cref{lem:balls}(b) for $D'=r$ and $\delta=\gamma/4$, and let $d_0=\max(d_1,2)$.
Suppose that (i) and~(ii) hold at some $d\ge d_0$, and let $X$ be the set of output positions at which the merge depth
exceeds~$s$. By \cref{thm:HCprime}(b) there is an index code with at most
\[
  \eps L+\frac{\beta(k+2)kC(L+1)}{\tau}+\eta L\le(\eps+\eta)L+\frac{\gamma(L+1)}4\le\Bigl(\frac12-\frac\gamma2\Bigr)L
\]
bits for a system $(I(\ell))_\ell$ with $I(\ell)\subseteq B^r(\ell)$. Some fibre of it has at least
$2^{(1/2+\gamma/2)L}$ elements, and it is a binary storage code for the relation $N(q)=I(\tpos(q))$
\cite[Remark~4.9(b)]{FourthPaper2026}, with $N(q)\subseteq B^r(\tpos(q))$. By \cref{lem:balls}(b) the fibre has at most
$2^{(1/2+\gamma/4)L}$ elements, a contradiction. For the last sentence, a lineage assignment for the empty genie is one
for every~$\Gamma_{\mathrm F}$, and $\Gamma_{\mathrm F}$-free paths are paths; so the hypothesis gives~(ii), with $\eta=0$,
for a cold genie at $(V,\frac14)$, which exists by \cref{lem:cold}.
Formally, \lean{MergeDepth.cor_HCstar_step} is the corollary at one size: for a refinement of width~$m$ with the
hypotheses of \lean{MergeDepth.hop_charging_fields}, $T\le C(L+1)$, at most $\eta L$ exceptional positions and
$4\beta mkC\le\gamma\tau$, and with the bound $2^{(1/2+\gamma/4)L}$ on binary storage codes for relations inside the balls
of radius $(s+1)\tau D$ as a hypothesis, the hypotheses are contradictory. The large fibre is
\lean{MergeDepth.exists_large_fibre_card}, and \lean{MergeDepth.isStorageCode_fibre_of_indexCode} shows that it is a
storage code for $q\mapsto I(\tpos(q))$.
\end{outline}

Corollary~HC$^*$ is the deep-nesting companion of Corollary~F$'$~\cite[Corollary~6.3]{BoundedNesting2026}, which it
contains, since bounded entry leaves no deep merges. Given Question~9.1, an oblivious linear-time transposer must
therefore, whatever its code, merge deep lineages more than $s$ times along some path from each of more than $\eta L$
outputs, for all fixed~$s$ and $\eta<\frac12-\eps$ and all large~$d$. The proof uses Question~9.1 only through
\cref{lem:balls}(b). An all-code crux for ball systems would serve as well, since the index code of
\cref{thm:HCprime}(b) is also one for the ball system: $\Crux(\Sball_{k,C,\beta,r},\rho)$, in the sense
of~\cite[Definition~4.8]{FourthPaper2026} for the class of~\cite[Definition~7.4]{BoundedNesting2026}, with
$\rho>\eps+\eta$ and $\beta(k+2)kC/\tau<(\rho-\eps-\eta)/4$, excludes the same machines. The crux for the cone class
$\mathcal S_{k,C,\beta,D}$ does not apply directly, because the systems of \cref{thm:HCprime}(b) are not cone systems of
whole-step genies of bounded full depth.

\subsection{Routing through deep hops, and the number of deep cells}\label{sec:hop-routing}

\begin{proposition}[Routing through deep hops is hop-limited]\label{prop:hop-routing}
\partformal{MergeDepth/HopRouting.lean}{MergeDepth.hop_routing_a,MergeDepth.hop_routing_b}{the bound $AL/\log_2d$ on the
number of positions~$\ell$ with $\tpos(\ell)\in B^{D''}(\ell)$, which is \cref{lem:balls}(a); the formal statements
count these positions}
Suppose that $d\ge2$ and that $M$ appends exactly $L$ output bits. Let $\Gamma_{\mathrm F}$ be a cold genie at $(V,\eps)$
with depth~$D$, let $D''\ge D$ be an integer, and put
\[
  K''=\sum_{i=0}^{D''}(k+1)^i,\qquad A=4C(K''+1)(k+2)(k+6)+3,\qquad h=\Bigl\lceil\frac{D''+2-D}D\Bigr\rceil.
\]
\begin{enumerate}[label=(\alph*)]
\item For all but at most $AL/\log_2d$ output positions~$\ell$, every $\Gamma_{\mathrm F}$-free path of~$H_M$ from
$t_\ell$ to a step that reads position $\tpos(\ell)$ contains at least $h$ deep edges.
\item Let $c_0\ge0$, and let nonnegative weights $f_\ell(\pi)$ be given on the $\Gamma_{\mathrm F}$-free paths~$\pi$
of~$H_M$ from~$t_\ell$ to steps that read $\tpos(\ell)$, such that $\sum_\pi f_\ell(\pi)\le1$ for every~$\ell$ and
$\sum_\ell\sum_{\pi\ni e}f_\ell(\pi)\le c_0$ for every deep edge~$e$. Then
$\sum_\ell\sum_\pi f_\ell(\pi)\le AL/\log_2d+c_0N_{\mathcal R}/h\le AL/\log_2d+c_0kC(L+1)/h$.
\end{enumerate}
\end{proposition}

\begin{outline}
(a) Apply \cref{lem:balls}(a) to the relation $N(q)=B^{D''}(\tpos(q))$: at most $AL/\log_2d$ positions~$p$ satisfy
$p\in B^{D''}(\tpos(p))$. For every other output position~$\ell$, the position $p=\tpos(\ell)$ satisfies $\tpos(p)=\ell$
and $p\notin B^{D''}(\ell)$, so every path of~$H_M$ from~$t_\ell$ to a step that reads~$p$ has more than $D''$ edges. A
$\Gamma_{\mathrm F}$-free path with $h'$ deep edges splits at them into $h'+1$ $\Gamma_{\mathrm F}$-free paths without
deep edges, each with at most $D-1$ edges by \cref{lem:cold}, so it has at most $(h'+1)D-1$ edges. Hence
$(h'+1)D-1\ge D''+1$, that is, $h'\ge(D''+2-D)/D$, and $h'\ge h$.
(b) The exceptional positions of~(a) contribute at most $AL/\log_2d$. For the others every path carries at least $h$
deep edges, so $h$ times their total weight is at most $\sum_e\sum_\ell\sum_{\pi\ni e}f_\ell(\pi)$, summed over the deep
edges, which is at most $c_0N_{\mathcal R}$ by~\eqref{eq:hop-NR}.
Formally, \lean{MergeDepth.length_le_of_coldDep} is the splitting at deep edges, for a cold genie in the sense of
\lean{MergeDepth.IsColdGenie}, and \lean{MergeDepth.hop_routing_a} is~(a) for one position, with the ball
\lean{MergeDepth.stepBall}. Part~(b) is \lean{MergeDepth.hop_routing_b}, for weights on any finite sets of such paths and
the positions $\ell$ with $\tpos(\ell)\in B^{D''}(\ell)$ as the exceptional set; it rests on the double counting
\lean{MergeDepth.routing_count} and on \lean{MergeDepth.rCount_le_card_edgesOf}: the edges of a path of~$H_M$ are
distinct, since steps decrease along it.
\end{outline}

The proposition is unconditional. Since $A=O((k+1)^{D''})$, the exceptional fraction tends to zero for every fixed~$D''$,
and also when $D''$ grows with~$d$ so slowly that $(k+1)^{D''}=o(\log d)$; then $h$ grows too. Read heuristically: after
the cold genie, the deep hops are the only long-range carriers, routing through them delivers $o(L)$ once $h\to\infty$,
and whatever a transposer delivers beyond the $\eps L$ revealed bits needs a network code on the deep hops whose gain over
routing grows with $D''/D$.

\begin{proposition}[The number of deep cells]\label{prop:capacity}
\formal{MergeDepth/Capacity.lean}{MergeDepth.card_deep_mul_le,MergeDepth.card_deepHops_le,MergeDepth.two_mul_card_deepHops_le,MergeDepth.card_deep_le_of_lt}
We have $\card{\mathcal R}\le kT/(V+1)$ and $N_{\mathcal R}\le\min\{kT,\card{\mathcal R}(T-1)/2\}$. In particular,
$\card{\mathcal R}(T-1)/2<kT$ only if $\card{\mathcal R}<2kT/(T-1)$, which for $T>2k+1$ means $\card{\mathcal R}\le2k$.
\end{proposition}

\begin{outline}
Each head makes at most $T$ entries, and each pair of~$\mathcal R$ takes more than~$V$ of them; this gives the first
bound, as in~\cite[Section~2.2]{BoundedNesting2026}. The bound $kT$ is~\eqref{eq:hop-NR}. For $(h,c)\in\mathcal R$, the
hops of head~$h$ at~$c$ partition the interval from the first to the last visit of~$c$, of length at most $T-1$, into
half-open intervals $(p,u]$~\cite[Section~2.2, (P1)]{BoundedNesting2026}; the long ones have length at least two, so
there are at most $(T-1)/2$ of them. Summing over~$\mathcal R$ gives the second bound. Finally,
$2kT/(T-1)=2k+2k/(T-1)<2k+1$ when $T>2k+1$.
Formally, $\mathcal R$ is the eighth paper's \lean{BoundedNesting.heavyPairs} and the deep hops are the long hops outside its
\lean{BoundedNesting.lightHops} (\lean{MergeDepth.deepHops}); the first bound is its
\lean{BoundedNesting.card_heavyPairs_mul_le}, and the disjointness of the intervals is its hop uniqueness.
\end{outline}

Theorems~HC and~HC$'$, Corollaries~HC and~HC$^*$, \cref{prop:hop-routing} and \cref{prop:merge-reduction} below depend on
the deep cells only through~$N_{\mathcal R}$ and the code; none uses $\card{\mathcal R}$. A scratch of fewer than $\eps L$
cells, as in~\cite[Problem~10.5]{BoundedNesting2026}, gives them nothing beyond \cref{prop:capacity}, which improves the
hop count only when there are at most $2k$ deep cells.

\begin{remark}[Heuristic: capacity bounds storage time, not depth]\label{rem:capacity-heuristic}
This is a heuristic. The number of deep cells bounds the total length $\card{\mathcal R}T$ of the deep hops, that is, how
long the deep part holds values; but cold storage can carry the long-time transport, while the deep part supplies depth
through hops that may be short. Demand paths of the form ``cold jump, a few short deep hops, cold jump, \dots'' use deep
cell-time of order~$h$ each, so capacity allows $\card{\mathcal R}T/O(h)$ of them, the hop count only $O(kT/h)$. We
expect no proof to profit from a hypothesis $\card{\mathcal R}\le\delta L$ alone; the results above show only that none
of them uses it.
\end{remark}

\subsection{The merge DAG and code-level depth reduction}\label{sec:hop-merge}

After \cref{lem:short-paths}, every long path from an output alternates stretches of fewer than $\tau$ deep lineage hops
with deep merges. What remains is the pattern of the merges.

\begin{definition}[The merge DAG]\label{def:merge-dag}
\formal{MergeDepth/MergeDag.lean}{MergeDepth.MEdge,MergeDepth.MInit,MergeDepth.MCut}
In the setting of \cref{lem:short-paths}, let $Q$ be the set given there. The \emph{merge DAG}~$\mathcal M$ has as
vertices the nodes of~$G$ that are deep merges of $(\Gamma_{\mathrm F},Q)$-free paths starting at unrevealed output
nodes, and an edge $f\to f'$ whenever $f$ and~$f'$ are consecutive deep merges of such a path. A vertex is
\emph{initial} if it is the first deep merge of such a path.
\end{definition}

Steps decrease along paths, so $\mathcal M$ is acyclic. It depends on the code, that is, on the representation or the
pruning and on the lineage assignment, and on $\Gamma_{\mathrm F}$ and~$\tau$; beyond that it does not depend on the
schedule.

\begin{proposition}[Code-level depth reduction suffices]\label{prop:merge-reduction}
\partformal{MergeDepth/MergeSteps.lean}{MergeDepth.merge_reduction_machine}{for field DAGs, that the field DAG of a pruning
is a function DAG that refines the dependency DAG, as in \cref{thm:HCprime}(b); for such refinements, the last claim is
\lean{MergeDepth.merge_reduction_fields}}
In the setting of \cref{lem:short-paths}, with $G$ the coordinate DAG of an $\F$-linear data path of width~$m$ or the
field DAG of a pruning with $m=k+2$, let $s\ge0$, $r=(s+1)\tau D$, and let $Z$ be a set of vertices of~$\mathcal M$
that meets every path of~$\mathcal M$ with $s+1$ vertices whose first vertex is initial. Then every
$(\Gamma_{\mathrm F},Q\cup Z)$-free path of~$G$ from an output node that is not revealed for
$(\Gamma_{\mathrm F},Q\cup Z)$ has fewer than $r$ edges. If $M$ transposes all $d\times d$ matrices, the conclusion of
\cref{thm:HCprime}(a) holds with $X=\emptyset$ and the rank bound increased by $\card Z$, and that of
\cref{thm:HCprime}(b) with $X=\emptyset$ and the bit count increased by $\beta\card Z$.
\end{proposition}

\begin{outline}
Let $\pi$ be such a path, from an output node~$u_0$. It is $(\Gamma_{\mathrm F},Q)$-free and starts at an unrevealed
output node, so its deep merges $f_1,\dots,f_{s'}$ form a path of~$\mathcal M$ whose first vertex is initial. If
$s'\ge s+1$, then $Z$ contains one of $f_1,\dots,f_{s+1}$; it is a vertex of~$\pi$ other than~$u_0$, which is not in~$Z$,
so $\pi$ would not be $Z$-free after its first vertex. Hence $s'\le s$, and $\pi$ has fewer than $r$ edges by
\cref{lem:short-paths}. Apply \cref{lem:expansion} with $\Gamma=\Gamma_{\mathrm F}$, the set $Q\cup Z$ and $X=\emptyset$,
as in the proof of \cref{thm:HCprime}; a node of~$Z$ costs rank at most one in~(a) and at most $\beta$ bits in~(b).
Formally, the merge DAG and the cut condition are \lean{MergeDepth.MEdge}, \lean{MergeDepth.MInit} and
\lean{MergeDepth.MCut}, for any network, deep relation and lineage assignment. \lean{MergeDepth.dm_le_of_cut} bounds the
deep merges of the paths in question by~$s$, and \lean{MergeDepth.length_le_of_cut} adds the bound of
\lean{MergeDepth.length_le_of_cov} (\cref{lem:short-paths}). For a linear network, \lean{MergeDepth.merge_reduction}
uses \lean{MergeDepth.rigidity_core_rev}, the form of \cref{lem:expansion}(a) in which a revealed output node
contributes to~$R$ only; \lean{MergeDepth.merge_reduction_steps} passes to the coordinate DAG as in Theorem~HC$'$, with
the rank bound $\beta\card{\Gamma_{\mathrm F}}+\card Q+\card Z$, and \lean{MergeDepth.merge_reduction_machine} states
$\rank R\le\eps L+mkC(L+1)/\tau+\card Z$ with $\tau\card Q\le mkC(L+1)$.
\end{outline}

The proposition gives a sufficient condition only. It moves the remaining difficulty from the schedule to the code: on
any schedule, depth-robust or not, a code whose merge DAG has a cut of $o(L)$ vertices meeting all long merge paths is
excluded as in Corollaries~HC and~HC$^*$. Heuristically, a code that merges deep lineages at every deep visit has a merge
DAG as rich as the deep part of the schedule, and then no such cut need exist.

\begin{conjecture}[Conjecture HV$^*$: code-level depth reduction]\label{conj:HV}
Let $k,V\ge1$ and $\beta\ge2$ be integers, $C\ge1$ and $0<\eps\le1$. There is $c>0$ such that for every integer
$\tau\ge1$ there is an integer~$r$ with the following property. For every $\eta>0$ there is~$d_0$ such that, for every
$d\ge d_0$, every machine with $k$ work tapes and value width at most~$\beta$ that is oblivious at length $L=d^2b(d)$
with halting time at most $C(L+1)$ and transposes all $d\times d$ matrices has an index code with at most
$\eps L+cN_{\mathcal R}/\tau+\eta L$ bits for its ball system $(B^r(\ell))_\ell$, where $N_{\mathcal R}$ counts its deep
hops at~$V$. In the \emph{linear form}, for a field~$\F$ and an integer $m\ge1$, with $c$ allowed to depend on $m$
and~$\F$, the same holds for the machines that have an $\F$-linear data path of width at most~$m$ at length~$L$, with the
minrank over~$\F$ of the ball system in place of the length of an index code.
\end{conjecture}

\begin{remark}[On Conjecture HV$^*$]\label{rem:HV}
(a) \Cref{thm:HCprime} proves the conjecture, with $\eta=0$, $r=(s+1)\tau D$ and $c=\beta(k+2)$, or $c=m$ in the
linear form, for machines whose codes have merge depth at most~$s$ for some pruning or representation, cold genie and
lineage assignment; \cref{prop:merge-reduction} proves it when the merge DAG has a cut as there of at most $\eta L/\beta$
vertices. (b) Since $N_{\mathcal R}\le kC(L+1)$, the conjecture is equivalent to the statement that for every
$\eps'>0$ there is~$r$ such that these ball systems have index codes, or minrank, of at most $\eps'L+o(L)$: take
$\eps=\eps'/2$ and $\tau\ge4ckC/\eps'$ in one direction, any~$c$ in the other. The form with $N_{\mathcal R}$ records the
rate at which~(a) charges deep hops; $V$ does not affect the truth of the conjecture. (c) The conjecture and an
affirmative answer to Question~9.1 for every degree bound exclude every oblivious linear-time transposer: for a
machine~$M$ take $V=1$, $\beta=\max(2,\beta(M))$, $\eps=\frac18$, $\tau\ge8ckC$ and $\eta=\frac1{16}$; at large~$d$ the
ball system has an index code with at most $\frac7{16}L$ bits, so some fibre, a storage code for the relation
$B^r(\tpos(q))$, has at least $2^{9L/16}$ elements, contrary to \cref{lem:balls}(b) with $\delta=\frac1{32}$. Likewise the
linear form and $\LinCrux_\F(\Sball_{k,C,\beta,r},\rho)$, at the radius~$r$ attached to some $\tau\ge8ckC/(\rho-\eps)$
with $\eps<\rho$, exclude linear ones, as in the proof of Corollary~HC. (d) Heuristically, the conjecture is of
network-coding type: after the cold genie the deep hops are the only long-range carriers (\cref{prop:hop-routing}), and
the conjecture bounds what they deliver against side information in balls of bounded radius by a constant times their
number, divided by~$\tau$. In general networks such bounds fail: coding can beat routing by a factor that grows with the
size of an undirected network~\cite{BravermanHe2026}. So the conjecture can hold only because of the shape of tape
schedules and the transposition demand, and a counterexample would be a tape-shaped network code on deep hops that merges
deep lineages along a pattern without cheap cuts, with a gain over routing that grows with the radius: the mechanism a
construction needs.
\end{remark}
